\section{Introduction}
Deep classifiers can change their predictions under small adversarial
perturbations \cite{szegedy2014intriguing,goodfellow2015explaining}.
Adversarial training addresses this failure by optimizing against perturbed
examples and remains a strong empirical defense \cite{madry2018towards}. Its
main obstacle is cost because multi-step training attacks require several
forward-backward passes for nearly every parameter update.

Fast adversarial training replaces the inner loop with a single FGSM step.
With randomized initialization and a suitable schedule, it can approach the
robustness of PGD training at much lower cost \cite{wong2020fast}. Yet
single-step training can exhibit \emph{catastrophic overfitting} (CO):
iterative-attack accuracy falls sharply toward zero while clean or FGSM
accuracy remains high. FastAdv+ responds using periodic PGD validation
\cite{li2020towards}; GradAlign regularizes local input-gradient geometry at
every update \cite{andriushchenko2020understanding}. These approaches improve
reliability but add attack or gradient computation.

We ask whether the FGSM update itself provides an early warning. It computes an
input gradient at the attack starting point, and the endpoint input gradient is
obtained from the same parameter-backward pass without an additional forward or
backward pass. Their cosine similarity therefore requires no additional model
evaluation. We use ``attack-free monitor'' to mean that the monitor itself does
not run an auxiliary validation attack. The base training still generates FGSM
adversarial examples, and validation PGD-10 is still used to select
\emph{best\_robust.pt}. When the cosine-drop score crosses a frozen threshold,
the primary \emph{immediate-switch} response retains the completed update,
switches the next mini-batch to PGD-2, and continues with PGD-2. The matched
\emph{immediate-replay} ablation instead rolls back to the epoch-start state
before recovery.

The contributions are:
\begin{itemize}
    \item A gradient-consistency safety trigger calibrated on three CIFAR-10
    development seeds and audited on seven disjoint held-out events. We also
    test unchanged-threshold architecture transfer and protocol transfer after
    one CIFAR-100-specific calibration.
    \item An event-level timing study that records warning, unprotected CO,
    and actual PGD-2 start times. Immediate-switch and immediate-replay act
    before collapse in 5/5 matched events, compared with 3/5 for next-epoch
    escalation; five seeds do not establish a clear proportional advantage.
    \item Five-seed component controls and controlled FastAdv+/GradAlign
    reimplementations under one harness, with paired bootstrap intervals.
    \item A random-start stress test: the monitor stays silent on eight stable
    $8/255$ runs but detects 3/3 natural CO events at $16/255$; the tested
    immediate-replay ablation restores non-trivial robustness.
\end{itemize}

The study concerns empirical $\ell_\infty$ robustness. Certified defenses
provide formal guarantees through different mechanisms and are complementary
to the training-time monitor studied here.
